There are words in today's Gospel that the first servant ought to
recognize. Another man is at his feet, begging for patience. He has heard
this plea before. He has spoken it himself.

``Have patience with me, and I will pay thee all.'' Before the king, those
were his words. Now they come from his fellow servant. The voice has
changed, but the need is familiar. What will he do when he hears his own
prayer coming from someone else?

To understand his answer, we must first stand with him before the king.
He owes ten thousand talents. Jesus tells us that he has no means to pay.
Everything is threatened: his freedom, his family, everything he owns.
He falls down and asks for time. He promises to pay everything, although
he has nothing with which to begin.

The king gives him something much greater than time. He lets him go and
forgives the debt. There is no installment plan waiting outside. The man
has been released by compassion. Before he has done anything merciful for
another person, mercy has opened a future for him.

Saint John Chrysostom asks us to notice precisely that difference. The
servant requested a delay; the king gave complete forgiveness. If we
begin the story with what the neighbor owes, we miss its beginning. The
first fact about this man is the gift he could never have supplied for
himself.

That is also where our life before God begins. We did not create
ourselves. We did not send the Son into the world. We cannot turn pardon
into something God owes us because we have become sufficiently useful.
We receive life and redemption from him. Today's Collect calls us God's
family and asks for his continuing mercy. Even our good works belong to
a life that still needs his protection.

Then the servant goes out. The man who was helpless before the king is
suddenly powerful before another debtor. He takes him by the throat.
His demand is payment. His companion falls down and asks for patience.
The Gospel's answer is terribly short: ``And he would not.''

The smaller debt is real. Jesus has not said that nothing was owed.
But listen to the terrible difference between the two encounters. This
man received more than he asked for, and refuses to give even what is
asked of him. His freedom has become the means of taking away someone
else's freedom. He remembers his claim and forgets his rescue.

Saint Augustine brings the two encounters into one life. Each of us is
God's debtor; each of us also knows what it means for another person to
owe us something. We cannot divide those identities, becoming humble
petitioners before God and merciless creditors everywhere else. The
person asking for patience may stand where we ourselves have stood.

This is why the ending is so severe. The king calls the servant back
and reminds him of the whole debt forgiven. Then he hands him over to
punishment. Jesus directs the warning to his disciples: forgiveness must
come from the heart. We cannot set that ending aside and keep only the
comforting first release.

From the heart: that reaches deeper than a courteous answer. Augustine
warns that a gentle voice can conceal a relentless heart. We might say
that everything is forgiven and continue to desire another's misery.
Jesus asks for the surrender of that desire. The king's question is
personal: having received compassion, should this servant not have shown
compassion too? He could not pay the enormous debt. He could have loosened
his hands. Mercy asks him to do what lies before him with the freedom
he has been given.

Nor should we conclude that God's forgiveness was unreal. Saint Thomas
Aquinas makes that distinction carefully. The earlier pardon really
happened. A later sin does not undo it. But a new mortal sin can
turn us from God again, and contempt for mercy makes the new offense
more grievous. The gift was real. The refusal is real. What we do after
being forgiven matters for eternity.

That warning concerns the will's refusal of mercy. A wound that still
hurts is not the same thing. Today's Offertory names Job upright before
it tells of his terrible losses and bodily affliction. Pain does not
prove that the sufferer is the wicked servant. Nor does forgiveness
require calling an injury harmless. Augustine insists that inward
pardon can coexist with correction directed toward another's good.
We may tell the truth, seek restitution, and protect someone from further
harm while refusing to make hatred our purpose.

That distinction is where the Epistle becomes very practical. Saint
Paul tells us to put on God's armor, and then tells us whom we are
fighting: our struggle is not against flesh and blood. The devil seeks
our defeat. Our fellow servant must not become the enemy we live to
destroy. We can resist an unjust act without giving our own heart to
revenge.

Look at the armor Paul names. Truth, justice, faith, salvation, the
word of God; and on our feet, readiness for the Gospel of peace.
There is no protection here that requires us to become cruel. We do
not overcome spiritual wickedness by adopting its malice. The forgiven
servant's grip on another man's throat is a defeat, however powerful
he looks.

We may never seize anyone by the throat. But we can still choose to
make another person's humiliation the thing we want. We can keep an
old wrong ready to use whenever it will hurt most. We can refuse a
fair hearing because we have already decided that this person deserves
none. The Gospel asks us to examine that deliberate tightening of
our grip.

Bring one such encounter before God today. Begin with the mercy you
have needed yourself. Then ask what truth and charity require in this
particular relationship. Perhaps there is a patient hearing you can
give, a contemptuous remark you can leave unsaid, or a correction you
can make without trying to wound. Choose one actual act of mercy.
It need not pretend that trust has already been restored. It must
refuse to make another person's ruin your satisfaction.

The Mass gives us more than a command to try harder. At the offering,
the Secret asks for reconciliation and restored salvation through the
sacrifice God himself has willed. Christ's saving gift comes first.
At Communion, the Church can still ask God to judge persecution and
give help. Mercy does not silence the cry of the injured. It entrusts
judgment to God while resisting the desire to become an instrument
of cruelty ourselves.

Then the Postcommunion brings the whole question close to us. It
calls the Eucharist the food of immortality and asks for an inward
fruit of what the mouth has received. Christ gives himself to nourish
a life that reaches beyond death. His gift lays claim to our mind,
our heart, our will. The mercy we receive is to shape the person
who leaves the church and meets a neighbor.

The servant went out from the king forgiven. Outside, he heard his
own plea and would not listen. We too will go out. We will meet
people who owe us something, people whose faults are real, people
who need patience. We need not meet them empty of help. God's family
is guarded by continuing mercy, armed for peace, and fed by Christ.
Let us remember the gift when we hear the plea. Let the freedom
we have received become room for another person to receive mercy.
